% GENERATED FILE. Edit canonical lessons, then run npm run generate:books.
% canonical-source-hash: fnv1a64:fccd767090f237f0
% canonical-lessons: MR-C372-samsodhan, MR-C372-prayog, MR-C372-vidvan, MR-C372-prabandh, MR-C372-siddhant

\chapter{Research}
\label{ch:mr-research}
\hlchaptermodality{372}

\begin{chapteropening}
\textbf{By the end of this chapter:} \emph{I can talk about research, experiments and theories.}

Complete the last lesson of chapter 372: I can talk about research, experiments and theories.
\end{chapteropening}

\section[\texorpdfstring{\mr{संशोधन} (saṁśodhan)}{संशोधन (saṁśodhan)}]{\mr{संशोधन} (saṁśodhan) — research}

\label{lesson:MR-C372-samsodhan}



\begin{quote}
Before the new one: say the Marathi for to cram, then the Marathi for a certificate.
\end{quote}

\subsection*{You'll want to know: \mr{संशोधन}}
\textbf{\mr{संशोधन}} — \emph{saṁśodhan} — \textquotedblleft{}research\textquotedblright{}.

\begin{quote}
\mr{ती} \mr{पाण्याच्या} \mr{प्रदूषणावर} \mr{संशोधन} \mr{करते}.
\end{quote}

\begin{quote}
\emph{tī pāṇyācyā pradūṣṇāvar saṁśodhan karte.}
\end{quote}

\begin{quote}
She does research on water pollution.
\end{quote}

Neuter.

\mr{संशोधन} \mr{करणे}, to do research; the researcher is \mr{संशोधक}.

\begin{cousinweb}
\textbf{\mr{संशोधन}} (\emph{saṁśodhan}) is a tatsama, taken from Sanskrit \emph{saṁśodhana}, \textquotedblleft{}purifying; correcting\textquotedblright{} (Monier-Williams).
\end{cousinweb}

\begin{grammarlens}[title={What you've built}]
One more word for giving reasons and opinions.
\end{grammarlens}

\subsection*{Your turn}
\begin{itemize}
\raggedright
  \item \emph{Say it:} \emph{saṁśodhan}
  \item \emph{Say it:} \emph{saṁśodhan}, once more
  \item \emph{Say it:} \emph{tī pāṇyācyā pradūṣṇāvar saṁśodhan karte.} — She does research on water pollution.
  \item \emph{Recall:} say the Marathi for to cram, then the Marathi for a certificate, then say \emph{saṁśodhan} again
\end{itemize}

\begin{tcolorbox}[breakable,colback=teal!4,colframe=teal!35!black,title={Before you move on}]
What does \mr{संशोधन} mean? (\textquotedblleft{}Research\textquotedblright{}.) Say it once more.
\end{tcolorbox}
\section[\texorpdfstring{\mr{प्रयोग} (prayog)}{प्रयोग (prayog)}]{\mr{प्रयोग} (prayog) — an experiment}

\label{lesson:MR-C372-prayog}



\begin{quote}
Before the new one: say the Marathi for a certificate, then the Marathi for research.
\end{quote}

\subsection*{You'll want to know: \mr{प्रयोग}}
\textbf{\mr{प्रयोग}} — \emph{prayog} — \textquotedblleft{}an experiment\textquotedblright{}.

\begin{quote}
\mr{विद्यार्थ्यांनी} \mr{वर्गात} \mr{एक} \mr{सोपा} \mr{प्रयोग} \mr{केला}.
\end{quote}

\begin{quote}
\emph{vidyārthyānnī vargāt ek sopā prayog kelā.}
\end{quote}

\begin{quote}
The students did a simple experiment in class.
\end{quote}

Masculine.

\mr{प्रयोग} \mr{करणे}, to experiment; also a stage performance.

\begin{cousinweb}
\textbf{\mr{प्रयोग}} (\emph{prayog}) is a tatsama, taken from Sanskrit \emph{prayoga}, \textquotedblleft{}joining; application, use\textquotedblright{} (Monier-Williams).
\end{cousinweb}

\begin{grammarlens}[title={What you've built}]
One more word for giving reasons and opinions.
\end{grammarlens}

\subsection*{Your turn}
\begin{itemize}
\raggedright
  \item \emph{Say it:} \emph{prayog}
  \item \emph{Say it:} \emph{prayog}, once more
  \item \emph{Say it:} \emph{vidyārthyānnī vargāt ek sopā prayog kelā.} — The students did a simple experiment in class.
  \item \emph{Recall:} say the Marathi for a certificate, then the Marathi for research, then say \emph{prayog} again
\end{itemize}

\begin{tcolorbox}[breakable,colback=teal!4,colframe=teal!35!black,title={Before you move on}]
What does \mr{प्रयोग} mean? (\textquotedblleft{}An experiment\textquotedblright{}.) Say it once more.
\end{tcolorbox}
\section[\texorpdfstring{\mr{विद्वान} (vidvān)}{विद्वान (vidvān)}]{\mr{विद्वान} (vidvān) — a scholar; a learned person}

\label{lesson:MR-C372-vidvan}



\begin{quote}
Before the new one: say the Marathi for research, then the Marathi for an experiment.
\end{quote}

\subsection*{You'll want to know: \mr{विद्वान}}
\textbf{\mr{विद्वान}} — \emph{vidvān} — \textquotedblleft{}a scholar; a learned person\textquotedblright{}.

\begin{quote}
\mr{आजोबा} \mr{संस्कृतचे} \mr{मोठे} \mr{विद्वान} \mr{होते}.
\end{quote}

\begin{quote}
\emph{ājobā saṁskṛtce moṭhe vidvān hote.}
\end{quote}

\begin{quote}
Grandfather was a great scholar of Sanskrit.
\end{quote}

Masculine.

Also an adjective, learned: \mr{विद्वान} \mr{माणूस}.

\begin{cousinweb}
\textbf{\mr{विद्वान}} (\emph{vidvān}) is a tatsama, taken from Sanskrit \emph{vidvas}, \textquotedblleft{}knowing, wise\textquotedblright{} (Monier-Williams). Its nominative is \emph{vidvān}.
\end{cousinweb}

\begin{grammarlens}[title={What you've built}]
One more word for giving reasons and opinions.
\end{grammarlens}

\subsection*{Your turn}
\begin{itemize}
\raggedright
  \item \emph{Say it:} \emph{vidvān}
  \item \emph{Say it:} \emph{vidvān}, once more
  \item \emph{Say it:} \emph{ājobā saṁskṛtce moṭhe vidvān hote.} — Grandfather was a great scholar of Sanskrit.
  \item \emph{Recall:} say the Marathi for research, then the Marathi for an experiment, then say \emph{vidvān} again
\end{itemize}

\begin{tcolorbox}[breakable,colback=teal!4,colframe=teal!35!black,title={Before you move on}]
What does \mr{विद्वान} mean? (\textquotedblleft{}A scholar; a learned person\textquotedblright{}.) Say it once more.
\end{tcolorbox}
\section[\texorpdfstring{\mr{प्रबंध} (prabandh)}{प्रबंध (prabandh)}]{\mr{प्रबंध} (prabandh) — a thesis; a dissertation}

\label{lesson:MR-C372-prabandh}



\begin{quote}
Before the new one: say the Marathi for an experiment, then the Marathi for a scholar.
\end{quote}

\subsection*{You'll want to know: \mr{प्रबंध}}
\textbf{\mr{प्रबंध}} — \emph{prabandh} — \textquotedblleft{}a thesis; a dissertation\textquotedblright{}.

\begin{quote}
\mr{तिने} \mr{इतिहासावर} \mr{प्रबंध} \mr{लिहिला}.
\end{quote}

\begin{quote}
\emph{tine itihāsāvar prabandh lihilā.}
\end{quote}

\begin{quote}
She wrote a thesis on history.
\end{quote}

Masculine.

\mr{प्रबंधाचा} \mr{विषय}, the topic of a thesis.

\begin{cousinweb}
\textbf{\mr{प्रबंध}} (\emph{prabandh}) is a tatsama, taken from Sanskrit \emph{prabandha}, \textquotedblleft{}band, connexion\textquotedblright{} (CDIAL 8695), later a connected composition.
\end{cousinweb}

\begin{grammarlens}[title={What you've built}]
One more word for giving reasons and opinions.
\end{grammarlens}

\subsection*{Your turn}
\begin{itemize}
\raggedright
  \item \emph{Say it:} \emph{prabandh}
  \item \emph{Say it:} \emph{prabandh}, once more
  \item \emph{Say it:} \emph{tine itihāsāvar prabandh lihilā.} — She wrote a thesis on history.
  \item \emph{Recall:} say the Marathi for an experiment, then the Marathi for a scholar, then say \emph{prabandh} again
\end{itemize}

\begin{tcolorbox}[breakable,colback=teal!4,colframe=teal!35!black,title={Before you move on}]
What does \mr{प्रबंध} mean? (\textquotedblleft{}A thesis; a dissertation\textquotedblright{}.) Say it once more.
\end{tcolorbox}
\section[\texorpdfstring{\mr{सिद्धांत} (siddhānt)}{सिद्धांत (siddhānt)}]{\mr{सिद्धांत} (siddhānt) — a theory; a doctrine}

\label{lesson:MR-C372-siddhant}



\begin{quote}
Before the new one: say the Marathi for a scholar, then the Marathi for a thesis.
\end{quote}

\subsection*{You'll want to know: \mr{सिद्धांत}}
\textbf{\mr{सिद्धांत}} — \emph{siddhānt} — \textquotedblleft{}a theory; a doctrine\textquotedblright{}.

\begin{quote}
\mr{हा} \mr{सिद्धांत} \mr{अजून} \mr{सिद्ध} \mr{झालेला} \mr{नाही}.
\end{quote}

\begin{quote}
\emph{hā siddhānt ajūn siddha jhālelā nāhī.}
\end{quote}

\begin{quote}
This theory has not been proved yet.
\end{quote}

Masculine.

\mr{सिद्ध}, proved, + \mr{अंत}, end: a settled conclusion.

\begin{cousinweb}
\textbf{\mr{सिद्धांत}} (\emph{siddhānt}) is a tatsama, taken from Sanskrit \emph{siddhānta}, \textquotedblleft{}established end, doctrine\textquotedblright{} (CDIAL 13404).
\end{cousinweb}

\begin{grammarlens}[title={What you've built}]
One more word for giving reasons and opinions.
\end{grammarlens}

\subsection*{Your turn}
\begin{itemize}
\raggedright
  \item \emph{Say it:} \emph{siddhānt}
  \item \emph{Say it:} \emph{siddhānt}, once more
  \item \emph{Say it:} \emph{hā siddhānt ajūn siddha jhālelā nāhī.} — This theory has not been proved yet.
  \item \emph{Recall:} say the Marathi for a scholar, then the Marathi for a thesis, then say \emph{siddhānt} again
\end{itemize}

\begin{tcolorbox}[breakable,colback=teal!4,colframe=teal!35!black,title={Before you move on}]
What does \mr{सिद्धांत} mean? (\textquotedblleft{}A theory; a doctrine\textquotedblright{}.) Say it once more.
\end{tcolorbox}
